\documentclass[10pt,a4paper]{amsart}
\usepackage[margin=19mm]{geometry}
\usepackage{amsmath,amssymb,mathtools}
\usepackage[scr=rsfs,cal=euler]{mathalfa}
\usepackage{xfrac}
\usepackage[unicode,hidelinks,bookmarks=false]{hyperref}
\newcommand{\ie}{i.e.\ }
\newcommand{\cf}{cf.\ }
\newcommand{\resp}{respectively\ }
\newcommand{\Subsection}{\subsection}
\providecommand{\nobreakdash}{\nobreak\hbox{-}}
\setlength{\emergencystretch}{2em}
\multlinegap0pt
\usepackage{amssymb}
\usepackage{xcolor}
\usepackage{enumerate}
\newtheorem{theorem}{Theorem}
\title{On isotropy group of locally finite derivations on $\mathbb{K}[X,Y]$}
\author{Luis Cid, Marcelo Veloso}
\date{Source: arXiv:2603.23709v1 (CC BY 4.0)}
\begin{document}
\maketitle

\noindent\emph{Source and licence.} On isotropy group of locally finite derivations on $\mathbb{K}[X,Y]$. Version arXiv:2603.23709v1 (CC BY 4.0), distributed by arXiv under the Creative Commons Attribution 4.0 International licence, http://creativecommons.org/licenses/by/4.0/, as recorded in the arXiv licence metadata for this exact version and shown on the abstract page https://arxiv.org/abs/2603.23709v1. Full licence notice: \texttt{licenses/arxiv-2603.23709-CC-BY-4.0.txt}.\par\smallskip

\noindent\textit{Editorial scope.} Counted unit: the article's theorem th3.8 (source0.tex lines 469--474). The article's complete original proof of it is delivered with it (source0.tex lines 476--532, one \texttt{proof} environment of 57 lines). No mathematics is written, repaired or re-derived: the statement and the proof are the article's own lines, in the article's own notation, and no retained line is substituted or re-flowed.  No further source-local context is needed: every cross-reference of every retained line resolves inside the two delivered spans.  Nothing was omitted from any retained span, so the package carries no omission record.  No retained line contains a citation, so the wrapper carries no bibliography and no bibliography entry.  No retained line contains a figure environment, an image-inclusion command or drawing code, so no image is delivered and the package declares no asset dependency.  Editorial formatting only, in addition: an AMS \texttt{amsart} preamble that restates the article's own declarations and macros used by the retained lines, namely the environments \texttt{theorem} on separate counters, together with the standard packages those lines need.  The retained environments print in this document as Theorem 1 in order of appearance, whereas the article prints the same items with the numbering of its own full document.  The complete native source of the article as distributed in the arXiv e-print for this exact version is bundled unchanged as \texttt{sources/197100/source0.tex}.
\begin{theorem} \label{th3.8}
		Let $D=(aX +Y)\dfrac{\partial}{\partial X}+aY\dfrac{\partial}{\partial Y}$ be a nonzero linear derivation of $\mathbb{K}[X,Y]$, where   $a \in \mathbb{K}^*$. Then 
        \[
        \operatorname{Aut}(\mathbb{K}[X,Y])_D= \{(\alpha X +\beta Y,\,\alpha Y ) \mid \alpha \in \mathbb{K}^*,\ \beta \in \mathbb{K}\}.
        \]
\end{theorem}

\begin{proof}
Assume that $\rho$ is an automorphism on $\mathbb{K}[X,Y]$ defined by
\[
\rho(X) = \sum_{i=0}^t a_i(X) Y^i, \qquad \rho(Y) = \sum_{j=0}^s b_j(X) Y^j,
\]
where  $a_i(X), b_j(X) \in \mathbb{K}[X]$. We want to determine when $\rho D = D \rho$. By definition, we have
\[
\rho(D(Y)) = a \rho(Y) = a \sum_{j=0}^s b_j(X) Y^j
\]
and
\begin{align*}
D(\rho(Y)) 
&= D\left( \sum_{j=0}^s b_j(X) Y^j \right)
=\sum_{j=0}^sab_j^{'}(X)XY^j+\sum_{j=0}^sb_j^{'}(X)Y^{j+1} + \sum_{j=1}^sjab_j(X) Y^j.
\end{align*}
Since $\rho(D(Y)) = D(\rho(Y))$, equating coefficients of each power of $Y$ yields:

\begin{equation*}
    \begin{cases} 
    b_0(X) = b_0'(X) X,\\
    b_s'(X) = 0, \\ 
   a b_j(X) =  a b_j'(X) X + b_{j-1}'(X) + ajb_j(X), \quad 0<j<s.\\
    \end{cases}
    \end{equation*}

The first equation $b_0(X)=Xb_0'(X)$ has general polynomial solution $b_0(X)=cX$ for some $c\in\mathbb{K}$. Substituting into the equation for $j=1$ (with $b_0'(X)=c$), we obtain  
\[
ab_1(X) = ab_1'(X)X + b_0'(X) + ab_1(X). 
\]
This implies  $0 = ab_1'(X)X + c$. Then $c=0$ and $b_1'(X)=0$, giving $b_0(X)=0$ and $b_1(X)=\alpha\in\mathbb{K}$. For $j\geq 2$, the equation $a(1-j)b_j(X) = ab_j'(X)X + b_{j-1}'(X)$ with $b_{j-1}'=0$ (established inductively) gives $a(1-j)b_j(X)=ab_j'(X)X$, i.e.\ $Xb_j'(X)=(1-j)b_j(X)$, whose only polynomial solution for $j\geq 2$ is $b_j(X)=0$. Thus
\[
\rho(Y) = \alpha Y, \quad \alpha \in \mathbb{K}^*.
\]

Now compute 
\[
\rho(D(X)) = \rho(aX + Y) = a \rho(X) + \alpha Y=\sum_{i=0}^t aa_i(X) Y^i+\alpha Y.
\]
On the other hand,
\begin{align*}
D(\rho(X)) 
&= \sum_{i=0}^t aa_i^{'}(X)XY^i+\sum_{i=0}^t a_i^{'}(X)Y^{i+1} + \sum_{i=1}^t iaa_i(X)Y^i.
\end{align*}
From $\rho(D(X))=D(\rho(X))$, equating coefficients for each power of $Y$ gives the system:
\begin{equation*} 
    \begin{cases}
    a_s'(X) = 0,\\
    a_0(X) = a_0'(X) X,\\
    \alpha = a a_1'(X) X + a_0'(X), \\ 
   a(1-i) a_i(X) = a a_i'(X) X + a_{i-1}'(X), \quad 1<i\leq s.\\
    \end{cases}
    \end{equation*}
Solving this system gives $a_0(X)=\gamma X$ and  $\alpha = a\cdot 0 + \gamma = \gamma$. So $a_0(X)=\alpha X$. From the equations for $i\geq 2$ one finds $a_i(X)=0$, and the  equation for $i=1$ we obtain $a_1(X)=\beta\in\mathbb{K}$. Therefore, 
\[
\rho=(\alpha X + \beta, \, \alpha Y ), \mbox{ where } \quad \alpha,\beta \in \mathbb{K} \mbox{ and }\ \alpha\neq 0.
\]
\end{proof}

\end{document}
